% ============================================================
%  SUPPLEMENT 02 — NORM GEOMETRY IN HIGH DIMENSIONS
%  L_p Norms, Normalization Distortion, and Certification Geometry
%  Companion to: ECDR_Data_Analysis.tex
% ============================================================
\documentclass[11pt, a4paper]{article}

\usepackage{fontspec}
\setmainfont{Noto Serif}
\setsansfont{Noto Sans}
\setmonofont{Noto Sans Mono}[Scale=0.88]

\usepackage[english]{babel}
\usepackage{amsmath, amssymb, amsthm}
\usepackage{mathtools}
\usepackage[margin=2.5cm]{geometry}
\usepackage{booktabs, tabularx, array}
\usepackage{xcolor}
\newcolumntype{Y}{>{\raggedright\arraybackslash}X}
\newcolumntype{P}[1]{>{\raggedright\arraybackslash}p{#1}}
\usepackage{hyperref,xurl}
\usepackage{enumitem}
\usepackage{microtype}
\usepackage{parskip}
\setlength{\emergencystretch}{3em}

\definecolor{ink}{RGB}{20,20,20}
\definecolor{teal}{RGB}{0,120,120}
\definecolor{muted}{RGB}{90,90,90}

\hypersetup{colorlinks=true, linkcolor=teal, urlcolor=teal, citecolor=teal}

\theoremstyle{definition}
\newtheorem{definition}{Definition}[section]
\newtheorem{example}{Example}[section]
\newtheorem{remark}{Remark}[section]
\theoremstyle{plain}
\newtheorem{theorem}{Theorem}[section]
\newtheorem{lemma}{Lemma}[section]
\newtheorem{proposition}{Proposition}[section]

\newcommand{\R}{\mathbb{R}}
\newcommand{\norm}[2]{\left\|#1\right\|_{#2}}
\newcommand{\Se}{\mathcal{S}_e}
\newcommand{\code}[1]{{\small\ttfamily\nolinkurl{#1}}}

\usepackage{titlesec}
\titleformat{\section}{\large\bfseries\color{teal}}{}{0em}{\thesection\quad}
\titleformat{\subsection}{\normalsize\bfseries\color{ink}}{}{0em}{\thesubsection\quad}

\begin{document}

\begin{center}
  {\LARGE\bfseries\color{teal} Norm Geometry in High Dimensions}\\[0.5em]
  {\large\color{muted} $L_p$ Norms, Normalization Distortion, and Certification Geometry in $\R^{768}$}\\[1.2em]
  {\normalsize Supplement~02 to \textit{ECDR Data Analysis} \textbar{} October~2026}\\[0.3em]
  {\small\color{muted} This document is supplementary only. The main document is not modified.}
\end{center}

\vspace{1em}\hrule\vspace{1em}

\tableofcontents
\clearpage

% ============================================================
\section{The $L_p$ Family of Norms}

For a vector $v \in \R^n$, the $L_p$ norm is defined as:
\begin{equation}
  \norm{v}{p} \;\coloneqq\; \left(\sum_{i=1}^n |v_i|^p\right)^{1/p}, \quad p \ge 1.
\end{equation}
Two special cases are critical for ECDR certification:

\begin{definition}[$L_2$ Euclidean Norm]
  \[
    \norm{v}{2} = \sqrt{\sum_{i=1}^n v_i^2}.
  \]
  This norm induces Euclidean distance. The closed unit ball is $\{v:\|v\|_2\le1\}$ and includes its interior; only the boundary sphere $\{v:\|v\|_2=1\}$ consists of points at distance 1 from the origin.
\end{definition}

\begin{definition}[$L_\infty$ Chebyshev Norm]
  \[
    \norm{v}{\infty} = \max_{i \in \{1,\ldots,n\}} |v_i|.
  \]
  The $L_\infty$ norm measures the \emph{largest coordinate deviation}.
  Its unit ball is a \emph{hypercube}: all points with all coordinates in $[-1, 1]$.
\end{definition}

\subsection{Why $L_\infty$ Is Preferred for Box Domains}

In ECDR, the perturbation domain is often a \textbf{hyperrectangle} (box):
\[
  \mathcal{B} = \prod_{i=1}^{768} [-b_i, b_i]
  = \bigl\{ v \in \R^{768} : \norm{v}{\infty} \le \bar{b} \bigr\}
  \quad \text{(when all $b_i = \bar{b}$).}
\]
The $L_\infty$ ball is exactly a box. IBP (Interval Bound Propagation) operates
natively with box inputs, propagating coordinate-wise intervals through each neural
network layer. This makes $L_\infty$ perturbation the natural choice for certifiers.

% ============================================================
\section{Norm Equivalence in $\R^n$}

\begin{theorem}[Standard Norm Bounds]
  For any $v \in \R^n$:
  \begin{equation}\label{eq:norm_bounds}
    \norm{v}{\infty}
    \;\le\; \norm{v}{2}
    \;\le\; \sqrt{n}\,\norm{v}{\infty}.
  \end{equation}
\end{theorem}

\begin{proof}
  \textbf{Left inequality:} $\norm{v}{\infty}^2 = \max_i v_i^2 \le \sum_i v_i^2 = \norm{v}{2}^2$.
  Take square roots.

  \textbf{Right inequality:} $\norm{v}{2}^2 = \sum_{i=1}^n v_i^2 \le n \max_i v_i^2 = n\,\norm{v}{\infty}^2$.
  Take square roots: $\norm{v}{2} \le \sqrt{n}\,\norm{v}{\infty}$.
\end{proof}

\subsection{The ECDR-Specific Case: $n = 768$}

The window tensor $Z \in \R^{6 \times 128}$ is vectorized to $\R^{768}$ (6 channels
$\times$ 128 time steps). The norm bounds become:
\begin{equation}
  d_\infty \;\le\; d_2 \;\le\; \sqrt{768}\,d_\infty \;\approx\; 27.7128\,d_\infty,
\end{equation}
where $d_\infty = \norm{\Delta z}{\infty}$ and $d_2 = \norm{\Delta z}{2}$ for
the difference $\Delta z = Z_{k+1} - Z_k$ between consecutive windows.

\subsection{What This Means for Reading the Data Table}

The main data analysis reports empirical distance percentiles in both norms:

\begin{center}
\small
\begin{tabular}{@{}lrrr@{}}
\toprule
Pair Subset / Norm & p50 & p90 & p99 \\
\midrule
Same-pure / $L_2$ & 4.936 & 46.42 & 65.26 \\
Same-pure / $L_\infty$ & 0.987 & 7.064 & 12.31 \\
\bottomrule
\end{tabular}
\end{center}

\begin{example}[Verifying the Bound]
  At the p50, $d_2 = 4.936$ and $d_\infty = 0.987$.
  Check: $d_\infty \le d_2$? $\;0.987 \le 4.936$. \checkmark \\
  Check: $d_2 \le 27.71 \, d_\infty$? $\;4.936 \le 27.71 \times 0.987 = 27.35$. \checkmark \\
  The ratio of these two medians is about5.0; it is not necessarily
  the median of the per-vector ratios, and it does not prove correlation.
\end{example}

\begin{remark}[Quantile bounds, not exact quantile conversion]
For the same population, percentile$p$ and quantile convention, the pointwise
inequalities imply $Q_p(d_\infty)\le Q_p(d_2)\le\sqrt{768}Q_p(d_\infty)$
by monotonicity of quantiles. They do not give an exact conversion, the
joint distribution or the distribution of ratios. The reverse assertion
$Q_p(d_\infty)\le Q_p(d_2)/\sqrt{768}$ need not hold. The table above is
historical normalized same-pure-pair evidence; its fitting population must
not be substituted for the current frozen normalization without rechecking.
\end{remark}

% ============================================================
\section{Geometric Distortion Under Normalization}

\subsection{The Diagonal Scaling Map}

Let the channel-wise standardization be defined by means $\mu_c$ and standard
deviations $\sigma_c$ for $c \in \{1, \ldots, 6\}$. Each sample is transformed as:
\[
  Z[c, i] = \frac{X[c, i] - \mu_c}{\sigma_c}.
\]
The perturbation $\Delta x = X_t - X^r$ in raw space maps to
$\Delta z = N(X_t) - N(X^r) = D^{-1} \Delta x$ in normalized space, where
\[
  D = \operatorname{diag}(\underbrace{\sigma_1, \ldots, \sigma_1}_{128 \text{ times}},
                           \underbrace{\sigma_2, \ldots, \sigma_2}_{128 \text{ times}},
                           \ldots,
                           \underbrace{\sigma_6, \ldots, \sigma_6}_{128 \text{ times}})
  \in \R^{768 \times 768}.
\]

\subsection{Isotropic Ball in Normalized Space $\to$ Ellipsoid in Raw Space}

An $L_2$ ball of radius $\epsilon$ in normalized space is:
\[
  \mathcal{B}_{\text{norm}} = \bigl\{ \Delta z \in \R^{768} : \norm{\Delta z}{2} \le \epsilon \bigr\}.
\]
Its preimage under $D^{-1}$ (i.e., the set of raw-space perturbations $\Delta x$
that map into the ball) is:
\[
  \mathcal{E} = \left\{ \Delta x \in \R^{6 \times 128} :
    \sum_{c=1}^6 \frac{1}{\sigma_c^2} \sum_{i=0}^{127} \Delta x[c,i]^2 \le \epsilon^2 \right\}.
\]
This is a \textbf{hyperellipsoid} in raw space with semi-axes of length
$\epsilon \sigma_c$ along each channel direction. Channels with larger $\sigma_c$
admit larger raw-space perturbations for the same normalized radius.

\begin{example}[Semi-Axis Comparison]
  Using historical fitted parameters as an illustrative substitution (not the frozen K=12 constants):
  \begin{center}
  \small
  \begin{tabular}{@{}lrr@{}}
  \toprule
  Channel & $\sigma_c$ & Raw semi-axis at $\epsilon = 1$ \\
  \midrule
  \code{acc_x} & 0.413195 & 0.413 $g$ \\
  \code{acc_y} & 0.402293 & 0.402 $g$ \\
  \code{acc_z} & 0.356522 & 0.357 $g$ \\
  \code{gyro_x} & 0.417901 & 0.418 rad/s \\
  \code{gyro_y} & 0.406792 & 0.407 rad/s \\
  \code{gyro_z} & 0.261610 & 0.262 rad/s \\
  \bottomrule
  \end{tabular}
  \end{center}
  A certified $\epsilon = 1$ ball in normalized space corresponds to very different
  raw perturbation tolerances across channels. \code{gyro_z} tolerates the smallest
  raw perturbation ($\approx 0.26$ rad/s) before exiting the certified set.
\end{example}

\subsection{Why Mixing Units Without Normalization Is Wrong}

Raw accelerometer values are in gravitational units ($g$), while gyroscope values
are in rad/s. These are physically incommensurable quantities: adding them as if
they were the same type of number violates dimensional analysis.

Specifically, the ``naive'' raw $L_2$ distance adds $g^2$ and $(\mathrm{rad/s})^2$:
\[
  d_{\text{naive}} = \sqrt{\sum_{c=1}^3 \sum_i \Delta x_{\text{acc}}[c,i]^2
                           + \sum_{c=4}^6 \sum_i \Delta x_{\text{gyro}}[c,i]^2},
\]
where the first sum has units $g^2$ and the second has units $(\mathrm{rad/s})^2$.
The root has \emph{no consistent unit}. Normalization resolves this by mapping both
into dimensionless standard-deviation units.

% ============================================================
\section{High-Dimensional Geometry: Intuitions That Fail}

\subsection{Almost All Volume Is Near the Shell}

For a fixed relative shell thickness, the fraction of ball volume in that shell increases with dimension: an $L_2$ ball of radius $r$ in $\R^n$ has most
of its volume concentrated in a thin shell near the surface.

Formally, for $\epsilon < r$:
\[
  \frac{\operatorname{Vol}\bigl(B(r) \setminus B(r - \epsilon)\bigr)}
       {\operatorname{Vol}\bigl(B(r)\bigr)}
  = 1 - \left(1 - \frac{\epsilon}{r}\right)^n \;\xrightarrow{n \to \infty}\; 1.
\]
In $\R^{768}$, with $\epsilon = 0.01r$, the fraction of volume outside the
innermost 1\% shell is $1 - 0.99^{768} \approx 1 - 0.00041 = 99.96\%$.

\textbf{Implication for ECDR:} A perturbation sampled uniformly by volume in the ball is likely to lie near its boundary in high dimension. Arbitrary temporal sensor distributions need not be uniform and can concentrate near the center. The margin
bound $\underline{m}_{cj}$ must hold at all points in $\Se$, including those on
the far shell.

\subsection{Coordinate-Wise Coverage Collapse}

Marginal coverage does not imply joint coverage. If 768 coordinate events
were independent with probability 0.9 each, their intersection probability
would be $0.9^{768}$, about $7.2\times10^{-36}$. Independence is a toy
assumption, not a proved HAPT property. Without dependence assumptions,
$\max(0,1-\sum_i\Pr(E_i^c))\le\Pr(\cap_iE_i)\le\min_i\Pr(E_i)$.
Correlated data can have high or low joint coverage; a low-dimensional
manifold and near-zero coverage cannot be concluded from marginal percentiles.

\subsection{Consequence: Compare Structured Domains Under Matched Conditions}
Structured domains may reduce support penalties when their directions align
with margin geometry. They can also miss queries or remain too loose to
certify. H10--H11 therefore require actual K=12 regional bounds and declared
matching conditions. Neither low-rank PCA nor a diagonal box guarantees
simultaneously useful coverage and positive margins. The normalized box is
the minimal baseline; compare wave/structured candidates rather than assuming
the conclusion. $C=6,W=128,K=12$ gives768 input coordinates regardless of K.

% ============================================================
\section{$L_\infty$ Box Domains and the Normalization Constraint}

\subsection{Box in Normalized vs.\ Raw Space}

An $L_\infty$ box in normalized space:
\[
  \mathcal{B}_\infty = \bigl\{ \Delta z \in \R^{768} : |\Delta z[c,i]| \le b \;\forall c, i \bigr\}
\]
corresponds to a box in raw space with channel-dependent half-widths $b \cdot \sigma_c$:
\[
  \mathcal{B}_{\text{raw}} = \bigl\{ \Delta x \in \R^{6 \times 128} : |\Delta x[c,i]| \le b \cdot \sigma_c \;\forall c, i \bigr\}.
\]

\subsection{The Non-Linearity Trap: Dynamic Per-Window Normalization}

If instead of a fixed $N$ using global parameters $(\mu_c, \sigma_c)$, one applied
per-window z-score normalization:
\[
  Z_w[c, i] = \frac{X_w[c, i] - \bar{X}_w[c]}{\hat{\sigma}_w[c]},
\]
then the map $N_w$ depends on $X_w$ itself. The frozen affine
transformation identities no longer apply unchanged. A sound verifier could
in principle bound nonlinear normalization with suitable operators, positivity
conditions on the denominator and numerical enclosures; IBP/LiRPA is not
restricted in principle to affine or piecewise-linear graphs. That would be a
different reference function and contract. This study preserves frozen
train-fit statistics, including their actual stored dtype and values.

\begin{remark}
  This is why the data analysis mandates that normalization parameters be
  ``frozen prior to evaluation'' and not re-fitted per window or per inference call.
\end{remark}

% ============================================================
\section{Summary of Key Relationships}

\begin{center}
\small
\begin{tabularx}{\linewidth}{@{}lYY@{}}
\toprule
Concept & Formula & ECDR Implication \\
\midrule
Norm bounds & $d_\infty \le d_2 \le \sqrt{768}\,d_\infty$ & Cannot translate percentile of one norm to the other. \\
Normalization & $\Delta z = D^{-1} \Delta x$ & $L_2$ ball in $\R^{768}$ $\leftrightarrow$ ellipsoid in raw space. \\
$L_\infty$ box & $|\Delta z[c,i]| \le b$ & Raw equivalent: $|\Delta x[c,i]| \le b \sigma_c$ per channel. \\
Volume shell & $1-(1-\epsilon/r)^{768}$ & Most of ball volume is near the surface in $\R^{768}$. \\
Non-linear norm & Dynamic $N_w$ & Changes reference; needs a separately justified transform bound. \\
\bottomrule
\end{tabularx}
\end{center}

\begin{thebibliography}{9}
\bibitem{main}
  ECDR Research Workspace. \textit{ECDR Data Analysis} (English edition).
  \code{ECDR_Data_Analysis.tex}, October~2026.
\bibitem{math}
  ECDR Research Workspace. \textit{ECDR Mathematics Concepts}.
  \code{ECDR_Mathematics_Concepts_EN.tex}, October~2026.
\end{thebibliography}

\end{document}
